\documentclass[11pt,oneside]{article}
\pdfminorversion=7
\usepackage{QuizStyle}
\hypersetup{pdftitle={20261117 Quiz for MATH301 Modern Algebra I},pdfauthor={Jungwoo Kim},pdfsubject={Quiz for MATH301 Modern Algebra I},pdfcreator={LaTeX}}
\begin{document}
\quiztitle{20261117 Quiz}

\begin{convention}
For the exercises below we follow Dummit--Foote's convention: rings and subrings need not have a multiplicative identity, and ring homomorphisms need not preserve identities. When a problem assumes an identity, that assumption is retained. All other notation follows MATH000.
\end{convention}

\begin{exercise}{7.1.3}
Let $R$ be a ring with identity and let $S$ be a subring of $R$ containing the identity. Prove that if $u$ is a unit in $S$, then $u$ is a unit in $R$. Give an example showing that the converse is false.
\end{exercise}
\begin{solution}
If $u\in S^\times$, there is $v\in S$ with $uv=vu=1_S=1_R$. Thus $u\in R^\times$. Conversely, take $S=\Z\subseteq R=\Q$. The element $2$ is a unit of $\Q$, but not of $\Z$.
\end{solution}

\begin{exercise}{7.1.5}
Decide which of the following sets are subrings of $\Q$:
\begin{exercisesubparts}
\item The rational numbers with odd denominators when written in lowest terms.
\item The rational numbers with even denominators when written in lowest terms.
\item The nonnegative rational numbers.
\item The squares of rational numbers.
\item The rational numbers with odd numerators when written in lowest terms.
\item The rational numbers with even numerators when written in lowest terms.
\end{exercisesubparts}
\end{exercise}
\begin{solution}
\begin{exercisesubparts}
\item \textbf{Yes.} Such numbers are precisely $a/b$ with $a,b\in\Z$ and $b$ odd. They are closed under differences and products, and contain $0$.
\item \textbf{No.} The set does not contain $0$ (whose reduced denominator is $1$).
\item \textbf{No.} It contains $1$ but not $-1$.
\item \textbf{No.} It contains $1$ but not $-1$.
\item \textbf{No.} It does not contain $0$.
\item \textbf{Yes.} A reduced rational has even numerator precisely when it can be expressed as $2a/b$ with $b$ odd. The expressions $2a/b$ are closed under subtraction and multiplication, and include $0$. This is a subring in the textbook's convention, although it does not contain $1$.
\end{exercisesubparts}
\end{solution}

\begin{exercise}{7.1.6}
Decide which of the following are subrings of the ring of all functions $[0,1]\to\R$, with pointwise operations:
\begin{exercisesubparts}
\item Functions $f$ such that $f(q)=0$ for every $q\in\Q\cap[0,1]$.
\item All polynomial functions.
\item Functions with only finitely many zeros, together with the zero function.
\item Functions with infinitely many zeros.
\item Functions $f$ such that $\lim_{x\to1^-}f(x)=0$.
\item All rational linear combinations of $\sin nx$ and $\cos mx$, where $m,n\in\{0,1,2,\ldots\}$.
\end{exercisesubparts}
\end{exercise}
\begin{solution}
\begin{exercisesubparts}
\item \textbf{Yes.} Vanishing at each rational point is preserved under subtraction and multiplication.
\item \textbf{Yes.} Polynomial functions are closed under both operations and contain $1$.
\item \textbf{No.} Choose an infinite proper subset $A\subset[0,1]$. Let $f=1$ and let $g$ be $-1$ on $A$ and $1$ outside $A$. Both have no zeros, but $f+g$ is neither zero everywhere nor nonzero at all but finitely many points.
\item \textbf{No.} Let $A\subset[0,1]$ have both $A$ and its complement infinite. Their characteristic functions $\mathbf1_A$ and $\mathbf1_{[0,1]\setminus A}$ each have infinitely many zeros, but their sum is $1$.
\item \textbf{Yes.} Sums, differences, and products of functions tending to $0$ as $x\to1^-$ also tend to $0$.
\item \textbf{Yes.} The product-to-sum identities express products of sines and cosines of integral frequencies as rational linear combinations of the same functions. The set is also closed under subtraction and includes $1=\cos0x$.
\end{exercisesubparts}
\end{solution}

\begin{exercise}{7.1.8}
Describe the center of the real Hamilton quaternions $\mathbb H$. Prove that $\{a+bi\mid a,b\in\R\}$ is a subring of $\mathbb H$ which is a field but is not contained in its center.
\end{exercise}
\begin{solution}
Write $q=a+bi+cj+dk$. If $q$ commutes with $i$, then $qi=iq$ gives $c=d=0$. Commuting also with $j$ then gives $b=0$. Thus $Z(\mathbb H)=\R$, and real scalars clearly commute with every quaternion. The subset $\R+\R i$ is closed under addition and multiplication; if $a+bi\ne0$, its inverse is $(a-bi)/(a^2+b^2)$, also in the subset. Hence it is a field, isomorphic to $\C$. It is not contained in the center because $ij=k\ne-k=ji$.
\end{solution}

\begin{exercise}{7.1.10}
Prove that if $D$ is a division ring, then the centralizer $C_D(a)=\{x\in D\mid xa=ax\}$ is a division ring for every $a\in D$ (cf. Exercise 7.1.9).
\end{exercise}
\begin{solution}
The centralizer contains $1$ and is closed under subtraction and multiplication. If $0\ne x\in C_D(a)$, then $xa=ax$ implies $a=x^{-1}ax$ and hence $ax^{-1}=x^{-1}a$. Thus $x^{-1}\in C_D(a)$, so every nonzero element of the centralizer is invertible there.
\end{solution}

\begin{exercise}{7.1.11}
Prove that if $R$ is an integral domain and $x^2=1$, then $x=\pm1$.
\end{exercise}
\begin{solution}
Since $(x-1)(x+1)=x^2-1=0$ and $R$ has no zero divisors, either $x=1$ or $x=-1$.
\end{solution}

\begin{exercise}{7.1.12}
Prove that any subring of a field which contains the identity is an integral domain.
\end{exercise}
\begin{solution}
Let $S\subseteq F$ be such a subring. Its identity is $1_F\ne0$, and multiplication is commutative. If $a,b\in S$ and $ab=0$, the absence of zero divisors in the field $F$ implies $a=0$ or $b=0$. Thus $S$ is an integral domain.
\end{solution}

\begin{exercise}{7.1.14}
Let $x$ be a nilpotent element of the commutative ring $R$ (cf. Exercise 7.1.13).
\begin{exercisesubparts}
\item Prove that $x$ is either zero or a zero divisor.
\item Prove that $rx$ is nilpotent for every $r\in R$.
\item Prove that $1+x$ is a unit in $R$.
\item Deduce that the sum of a nilpotent element and a unit is a unit.
\end{exercisesubparts}
\end{exercise}
\begin{solution}
Choose $m\ge1$ with $x^m=0$.
\begin{exercisesubparts}
\item If $x\ne0$, choose $m$ minimal. Then $m\ge2$ and $x^{m-1}\ne0$, while $x\cdot x^{m-1}=0$.
\item Commutativity gives $(rx)^m=r^mx^m=0$.
\item The finite geometric identity yields
\[
(1+x)\sum_{k=0}^{m-1}(-x)^k=1-(-x)^m=1.
\]
Hence $1+x$ is a unit.
\item If $u\in R^\times$, then $u+x=u(1+u^{-1}x)$. By (b), $u^{-1}x$ is nilpotent, so (c) shows that both factors are units.
\end{exercisesubparts}
\end{solution}

\begin{exercise}{7.1.21}
Let $X$ be a nonempty set, and let $\mathcal P(X)$ be its power set. Define $A+B=(A\setminus B)\cup(B\setminus A)$ and $A\times B=A\cap B$.
\begin{exercisesubparts}
\item Prove that $\mathcal P(X)$ is a ring under these operations.
\item Prove that this ring is commutative, has an identity, and is a Boolean ring.
\end{exercisesubparts}
\end{exercise}
\begin{solution}
For each $A\subseteq X$, let $\chi_A:X\to\Z/2\Z$ be its characteristic function. Pointwise,
\[
\chi_{A+B}=\chi_A+\chi_B,\qquad \chi_{A\times B}=\chi_A\chi_B.
\]
The map $A\mapsto\chi_A$ is a bijection from $\mathcal P(X)$ onto the ring of functions $X\to\Z/2\Z$. Hence (a) these operations satisfy all ring axioms, with zero $\varnothing$ and each set its own additive inverse. Moreover, (b) multiplication is commutative, its identity is $X$, and $A\times A=A$ for every $A$. Thus this is a Boolean ring.
\end{solution}

\Needspace{16\baselineskip}
\begin{exercise}{7.1.25}
Let $I$ be the ring of integral Hamilton quaternions and define the norm
\[
N:I\to\Z,\qquad N(a+bi+cj+dk)=a^2+b^2+c^2+d^2.
\]
\begin{exercisesubparts}
\item Prove $N(\alpha)=\alpha\overline\alpha$, where $\overline{a+bi+cj+dk}=a-bi-cj-dk$.
\item Prove $N(\alpha\beta)=N(\alpha)N(\beta)$ for all $\alpha,\beta\in I$.
\item Prove that $\alpha\in I$ is a unit if and only if $N(\alpha)=1$. Show that $I^\times\cong Q_8$.
\end{exercisesubparts}
\begin{hint}
For $\alpha\ne0$, its inverse in the rational quaternions is $\overline\alpha/N(\alpha)$.
\end{hint}
\end{exercise}
\begin{solution}
\begin{exercisesubparts}
\item Expanding with $i^2=j^2=k^2=-1$ and $ij=k=-ji$, $jk=i=-kj$, $ki=j=-ik$, all cross terms cancel in $\alpha\overline\alpha$, leaving $a^2+b^2+c^2+d^2$. This real scalar also equals $\overline\alpha\alpha$.
\item Quaternion conjugation reverses products: $\overline{\alpha\beta}=\overline\beta\,\overline\alpha$ (verified on $1,i,j,k$ and extended by bilinearity). Hence
\[
N(\alpha\beta)=\alpha\beta\overline\beta\,\overline\alpha
=N(\beta)\alpha\overline\alpha=N(\alpha)N(\beta).
\]
\item If $\alpha\beta=1$ in $I$, then $N(\alpha)N(\beta)=1$. Both norms are nonnegative integers, hence $N(\alpha)=1$. Conversely, if $N(\alpha)=1$, then $\overline\alpha\in I$ is its inverse. The integer solutions of $a^2+b^2+c^2+d^2=1$ are precisely $\pm1,\pm i,\pm j,\pm k$. Their products satisfy the standard quaternion-group relations, so $I^\times=\{\pm1,\pm i,\pm j,\pm k\}\cong Q_8$.
\end{exercisesubparts}
\end{solution}

\begin{exercise}{7.1.28}
Let $R$ be a ring with $1\ne0$. A nonzero $a$ is a \emph{left zero divisor} if $ax=0$ for some nonzero $x$, and a nonzero $b$ is a \emph{right zero divisor} if $yb=0$ for some nonzero $y$. An element $u$ has a left inverse if $su=1$ for some $s$, and has a right inverse if $ut=1$ for some $t$.
\begin{exercisesubparts}
\item Prove that $u$ is a unit if and only if it has both a right and a left inverse.
\item Prove that if $u$ has a right inverse, then $u$ is not a right zero divisor.
\item Prove that if $u$ has more than one right inverse, then $u$ is a left zero divisor.
\item Prove that in a finite ring every element with a right inverse is a unit.
\end{exercisesubparts}
\end{exercise}
\begin{solution}
\begin{exercisesubparts}
\item If $su=1=ut$, then $s=s(ut)=(su)t=t$, so $s$ is a two-sided inverse. The converse is immediate.
\item If $ut=1$ and $yu=0$, then $y=(yu)t=0$.
\item If $ut=ut'=1$ with $t\ne t'$, then $u(t-t')=0$ and $t-t'\ne0$. Also $u\ne0$.
\item If $ut=1$, the map $R\to R$, $r\mapsto ru$, is injective: $ru=su$ implies $r=(ru)t=(su)t=s$. Since $R$ is finite, this map is surjective; thus $su=1$ for some $s\in R$. Part (a) shows that $u$ is a unit.
\end{exercisesubparts}
\end{solution}

\Needspace{17\baselineskip}
\begin{exercise}{7.2.3}
Let $R$ be a commutative ring with identity. Define the ring $R[[x]]$ of formal power series to consist of all expressions $\sum_{n=0}^\infty a_nx^n$ with $a_n\in R$, with coefficientwise addition and multiplication
\[
\left(\sum_{n\ge0}a_nx^n\right)\left(\sum_{n\ge0}b_nx^n\right)
=\sum_{n\ge0}\left(\sum_{k=0}^n a_kb_{n-k}\right)x^n.
\]
Convergence is not considered.
\begin{exercisesubparts}
\item Prove that $R[[x]]$ is a commutative ring with $1$.
\item Show that $1-x$ is a unit with inverse $1+x+x^2+\cdots$.
\item Prove that $\sum_{n\ge0}a_nx^n$ is a unit if and only if $a_0$ is a unit in $R$.
\end{exercisesubparts}
\end{exercise}
\begin{solution}
\begin{exercisesubparts}
\item Addition gives an abelian group with zero series and coefficientwise inverses. Multiplication is commutative since $a_kb_{n-k}=b_{n-k}a_k$; its associativity follows by writing the coefficient of $x^n$ in either parenthesization as the finite sum $\sum_{i+j+k=n}a_ib_jc_k$. Distributivity holds coefficientwise, and the constant series $1_R$ is the identity.
\item For $g=\sum_{n\ge0}x^n$, the constant coefficient of $(1-x)g$ is $1$, and every coefficient of positive degree is $1-1=0$. Thus $(1-x)g=1$.
\item If $fg=1$, comparison of constant coefficients gives $a_0b_0=1$, so $a_0$ is a unit. Conversely, suppose $a_0\in R^\times$. Set $b_0=a_0^{-1}$ and recursively
\[
b_n=-a_0^{-1}\sum_{k=1}^na_kb_{n-k}\qquad(n\ge1).
\]
Then every coefficient of positive degree in $fg$, where $g=\sum b_nx^n$, is $0$, and its constant coefficient is $1$. Thus $fg=1$.
\end{exercisesubparts}
\end{solution}

\begin{exercise}{7.3.1}
Prove that the rings $2\Z$ and $3\Z$ are not isomorphic.
\end{exercise}
\begin{solution}
An additive-group isomorphism $\varphi:2\Z\to3\Z$ must send the generator $2$ to $3$ or $-3$. Since $4=2\cdot2=2+2$, multiplicative and additive preservation would give
\[
\varphi(4)=\varphi(2)^2=9,
\qquad\varphi(4)=2\varphi(2)=\pm6,
\]
a contradiction. Hence no ring isomorphism exists.
\end{solution}

\begin{exercise}{7.3.2}
Prove that the rings $\Z[x]$ and $\Q[x]$ are not isomorphic.
\end{exercise}
\begin{solution}
A ring isomorphism preserves units. The units of $\Z[x]$ are precisely $\{\pm1\}$, whereas the units of $\Q[x]$ are all nonzero constant rational polynomials, an infinite set. Therefore the two rings cannot be isomorphic.
\end{solution}

\begin{exercise}{7.3.3}
Find all homomorphic images of $\Z$.
\end{exercise}
\begin{solution}
If $\varphi:\Z\to S$ is a ring homomorphism, its additive kernel is an ideal $n\Z$ for a unique $n\ge0$. By the First Isomorphism Theorem,
\[
\im\varphi\cong\Z/n\Z.
\]
Conversely, for each $n\ge1$ the reduction map $\Z\to\Z/n\Z$ is surjective, and for $n=0$ the identity map has image $\Z=\Z/0\Z$. Thus these are exactly the homomorphic images, up to isomorphism (including the zero ring when $n=1$).
\end{solution}

\begin{exercise}{7.3.4}
Find all ring homomorphisms $\varphi:\Z\to\Z/30\Z$. In each case, describe the kernel and image.
\end{exercise}
\begin{solution}
Every additive homomorphism is determined by $e=\varphi(1)$ and has the form $\varphi_e(n)=[ne]_{30}$. It preserves multiplication precisely when $e^2=e$ modulo $30$. By the Chinese Remainder Theorem, $\Z/30\Z\cong\Z/2\Z\times\Z/3\Z\times\Z/5\Z$, and an idempotent has each coordinate $0$ or $1$. The eight choices and their kernels and images are:
\[
\begin{array}{@{}c c c@{}}
\toprule
 e & \ker\varphi_e & \im\varphi_e\\
\midrule
0&\Z&\{[0]_{30}\}\\
1&30\Z&\langle[1]_{30}\rangle\\
6&5\Z&\langle[6]_{30}\rangle\\
10&3\Z&\langle[10]_{30}\rangle\\
15&2\Z&\langle[15]_{30}\rangle\\
16&15\Z&\langle[16]_{30}\rangle\\
21&10\Z&\langle[21]_{30}\rangle\\
25&6\Z&\langle[25]_{30}\rangle\\
\bottomrule
\end{array}
\]
Here the image is the additive subgroup generated by $[e]_{30}$, and $\ker\varphi_e=(30/\gcd(30,e))\Z$ for $e\ne0$ (with $\ker\varphi_0=\Z$). Only $e=1$ preserves the ambient identity.
\end{solution}

\begin{exercise}{7.3.10}
Decide which of the following sets are ideals of $\Z[x]$:
\begin{exercisesubparts}
\item Polynomials whose constant term is divisible by $3$.
\item Polynomials whose coefficient of $x^2$ is divisible by $3$.
\item Polynomials whose constant, $x$-, and $x^2$-coefficients are all zero.
\item $\Z[x^2]$, the polynomials containing only even powers of $x$.
\item Polynomials whose coefficients sum to zero.
\item Polynomials $p(x)$ such that $p'(0)=0$.
\end{exercisesubparts}
\end{exercise}
\begin{solution}
\begin{exercisesubparts}
\item \textbf{Yes.} This is the kernel of evaluation at $0$ followed by reduction modulo $3$, $\Z[x]\to\Z/3\Z$.
\item \textbf{No.} The polynomial $1$ lies in the set but $x^2\cdot1=x^2$ does not.
\item \textbf{Yes.} This is the principal ideal $(x^3)$.
\item \textbf{No.} It contains $1$, but $x\cdot1=x$ is not in the set.
\item \textbf{Yes.} The sum of the coefficients of $p$ is $p(1)$, so this is the kernel of evaluation at $1$, namely $(x-1)$.
\item \textbf{No.} The polynomial $1$ satisfies $p'(0)=0$, whereas $x\cdot1=x$ does not.
\end{exercisesubparts}
\end{solution}
\end{document}
